\bookchapter{Chapitre 2 : Rangez votre chambre numérique — gestion des fichiers}

\chapterintro{Le smartphone devient encombré avec le temps ; de bonnes habitudes commencent par la gestion des fichiers. Comme votre pièce a besoin d'être rangée régulièrement, votre espace numérique a aussi besoin d'être entretenu. Un smartphone bien ordonné n'est pas seulement plus pratique pour retrouver les choses, il fonctionne aussi plus fluidement.}

\booksection{I. Système de stockage à trois niveaux : votre entrepôt numérique}

\begin{enumerate}
  \item \textbf{Stockage interne}: zone centrale, rapide, le système et les apps s'y trouvent. C'est l'espace le plus précieux de votre smartphone, à gérer avec parcimonie.
  \item \textbf{Carte SD (si prise en charge)}: placard d'extension, plus lent, adapté aux gros fichiers tels que films et photos. À noter : beaucoup de smartphones phares ne prennent plus en charge la carte SD.
  \item \textbf{Stockage cloud}: coffre-fort + disque réseau de synchronisation. Prévention contre la perte, permet aussi la synchronisation entre plusieurs appareils. iCloud, Baidu Netdisk, OneDrive sont des choix courants.
\end{enumerate}

\booksection{II. Connaître le « gestionnaire de fichiers »}

Android a généralement une app [Fichiers] ou [Gestionnaire de fichiers], iOS a l'app [Fichiers] intégrée. Ce sont vos « explorateurs » pour consulter tous les fichiers du smartphone.

\booksubsection{Signification des dossiers courants :}

\begin{itemize}
  \item \textbf{DCIM}: photos et vidéos prises par l'appareil photo.
  \item \textbf{Pictures}: autres images, telles que captures d'écran, images enregistrées.
  \item \textbf{Download}: fichiers téléchargés par le navigateur.
  \item \textbf{Music}: fichiers de musique.
  \item \textbf{Movies}: fichiers vidéo.
  \item \textbf{Documents}: fichiers de documents.
  \item \textbf{Android}: dossier système Android, contient les données des apps, à ne pas supprimer manuellement en général.
\end{itemize}

\booksection{III. Structure du système de fichiers : aller plus loin}

Android est basé sur Linux, le système de fichiers a une certaine structure hiérarchique. Bien que l'utilisateur ordinaire n'ait pas besoin de l'approfondir, connaître quelques concepts de base est utile.

\booksubsection{Partitions de stockage :}

\begin{itemize}
  \item \textbf{/system}: partition système, contient les fichiers du système d'exploitation. L'utilisateur ordinaire n'y a pas accès, nécessite des droits root.
  \item \textbf{/data}: partition de données, contient les données utilisateur et les données d'applications. C'est l'espace que vous utilisez au quotidien.
  \item \textbf{/sdcard ou /storage/emulated/0}: espace de stockage accessible à l'utilisateur, vos photos et fichiers téléchargés s'y trouvent.
\end{itemize}

\booksubsection{Stockage des données d'applications :}

\begin{itemize}
  \item \textbf{Stockage interne}: espace privé de l'app, inaccessible aux autres apps. Supprimé lors de la désinstallation de l'app.
  \item \textbf{Stockage externe}: fichiers stockés par l'app dans la zone accessible à l'utilisateur, tels que photos, téléchargements. Pas nécessairement supprimé lors de la désinstallation de l'app.
\end{itemize}

\booksection{IV. Sauvegarde et restauration de données : assurez vos données}

Les données n'ont pas de prix, la sauvegarde est une habitude à prendre.

\booksubsection{Stratégies de sauvegarde :}

\begin{enumerate}
  \item \textbf{Sauvegarde cloud}: automatique, pratique, mais nécessite du réseau et de l'espace cloud. iCloud, Huawei Cloud, Xiaomi Cloud, etc., tous prennent en charge.
  \item \textbf{Sauvegarde locale}: connectez un ordinateur avec un câble de données, sauvegardez avec l'outil officiel. Rapide, grand espace, mais opération manuelle.
  \item \textbf{Sauvegarde double}: sauvegarde cloud + sauvegarde locale, la plus sûre.
\end{enumerate}

\booksubsection{Que sauvegarder ?}

\begin{itemize}
  \item \textbf{Photos et vidéos}: prennent le plus de place, les plus précieuses.
  \item \textbf{Carnet de contacts}: les stocker dans un compte cloud est le plus pratique.
  \item \textbf{Historique de discussion WeChat}: utilisez la fonction de sauvegarde intégrée de WeChat.
  \item \textbf{Documents importants}: stockez-les sur un disque cloud ou un ordinateur.
\end{itemize}

\booksection{V. Sécurité des fichiers : protégez votre vie privée}

\begin{itemize}
  \item \textbf{Chiffrement des fichiers sensibles}: beaucoup de smartphones ont une fonction « dossier sécurisé » ou « espace privé », où vous pouvez placer des fichiers sensibles, accessibles uniquement par mot de passe ou empreinte.
  \item \textbf{Supprimer n'est pas détruire}: la suppression ordinaire marque seulement l'espace comme remplaçable, les données sont toujours là. Pour une suppression définitive, il faut utiliser un outil d'effacement sécurisé dédié, ou écrire par-dessus plusieurs fois.
  \item \textbf{Sécurité cloud}: choisissez un fournisseur cloud de bonne réputation, activez la double authentification, vérifiez régulièrement les appareils connectés.
\end{itemize}

\booksection{VI. Plus d'espace ? Faites un grand ménage}

\begin{itemize}
  \item \textbf{Vider le cache}: mémoire temporaire des apps, la vider libère de la place. Dans Réglages > Stockage.
  \item \textbf{Chercher les gros fichiers}: utilisez le gestionnaire de fichiers pour trier par taille, débusquez les « coupables » qui prennent de la place.
  \item \textbf{Supprimer les résidus}: après la désinstallation d'une app, il reste parfois des fichiers orphelins, à supprimer manuellement.
  \item \textbf{Nettoyer l'album}: supprimez les photos en double, les photos floues, les captures d'écran inutiles.
  \item \textbf{Nettoyer WeChat}: WeChat est un gros consommateur d'espace, nettoyer régulièrement l'historique de discussion et le cache est nécessaire.
\end{itemize}

\booksection{VII. Analyse complète des formats de fichiers courants}

Les extensions de fichiers sur le smartphone donnent le tournis ; les comprendre vous aide à mieux gérer vos fichiers :

\begin{center}
\begin{tabularx}{\linewidth}{|X|X|X|X|}
\hline
\textbf{Format} & \textbf{Type} & \textbf{Caractéristiques} & \textbf{Ouverture recommandée} \tabularnewline
\hline
JPG/JPEG & Image & Compression avec perte, petite taille, meilleure compatibilité. Format par défaut des photos du smartphone. & Album / n'importe quelle visionneuse d'images \tabularnewline
PNG & Image & Compression sans perte, prend en charge la transparence. Souvent utilisé pour captures d'écran, icônes. Taille plus grande que JPG. & Album / n'importe quelle visionneuse d'images \tabularnewline
MP3 & Audio & Compression avec perte, petite taille, excellente compatibilité. Le format audio le plus universel. & N'importe quel lecteur de musique \tabularnewline
FLAC & Audio & Compression sans perte, qualité sonore identique au CD, taille 3 à 5 fois plus grande que le MP3. Choix des audiophiles. & Lecteur de musique / app de lecture spécialisée \tabularnewline
WAV & Audio & Audio brut non compressé, meilleure qualité mais taille énorme. Courant en enregistrement professionnel. & Lecteur de musique \tabularnewline
M4A & Audio & Codage AAC, meilleure qualité que le MP3 à débit égal. Courant dans l'écosystème Apple. & Apple Music / lecteur de musique \tabularnewline
MP4 & Vidéo & Format conteneur vidéo le plus universel, excellente compatibilité. Format par défaut des vidéos enregistrées par le smartphone. & Album / n'importe quel lecteur vidéo \tabularnewline
MKV & Vidéo & Format conteneur open source, prend en charge plusieurs pistes audio et sous-titres. Courant pour les films. & VLC/nPlayer \tabularnewline
DOC/DOCX & Document & Document Word. DOCX est un format compressé, plus petit. DOC est l'ancien format. & WPS/Office \tabularnewline
PPT/PPTX & Présentation & Présentation PowerPoint. PPTX est un format compressé. & WPS/Office \tabularnewline
XLS/XLSX & Tableur & Tableur Excel. XLSX est un format compressé. & WPS/Office \tabularnewline
PDF & Document & Format multiplateforme, mise en page fixe non éditable. Courant pour contrats, mémoires, livres électroniques. & Adobe Reader/WPS \tabularnewline
ZIP & Compression & Format de compression le plus universel, bonne compatibilité multiplateforme. & Intégré au système / app de décompression \tabularnewline
RAR & Compression & Taux de compression supérieur à ZIP, mais nécessite un logiciel dédié pour décompresser. & App de décompression \tabularnewline
ISO & Image disque & Fichier image de disque, copie complète du contenu d'un disque. Courant pour les paquets d'installation système. & App de décompression / lecteur de disque virtuel \tabularnewline
HTML & Page web & Fichier de page web, ouvert avec un navigateur. Les pages web enregistrées utilisent souvent ce format. & Navigateur \tabularnewline
\hline
\end{tabularx}
\end{center}

\begin{tcbpractice}{Exercice : rangez le dossier « Téléchargements »}
Ouvrez le gestionnaire de fichiers > entrez dans le dossier « Téléchargements ». Regardez ce qu'il contient ? Supprimez l'inutile, classez l'utile (par exemple créez des dossiers « Documents professionnels », « Stickers » pour les y mettre). Profitez-en pour regarder l'espace de stockage, voyez quels types de fichiers prennent le plus de place.

\end{tcbpractice}

\booksection{VIII. Gestion des images et du système de fichiers : un patrimoine bien ordonné}

La section précédente vous a appris à « connaître les dossiers », mais ce n'est pas suffisant. Cette section traite\textbf{comment mettre en place activement une norme de gestion exécutable à long terme}, pour que vos fichiers soient toujours retrouvables, non perdus, non espionnés.

\booksubsection{Classification des fichiers et convention de nommage}

La source du chaos est l'absence de règles. Établir des règles est en fait simple, la clé est de\textbf{persévérer}.

\begin{itemize}
  \item \textbf{Méthode de nommage « date + événement »}: le nom du dossier contient un horodatage et une description d'événement, pour faciliter la mémoire et la recherche ultérieures. Par exemple\texttt{202606\_Tickets de voyage}, le\texttt{Bilan annuel 2025}, le\texttt{Anniversaire enfant 2024}. Vous voyez ainsi la chronologie et le contenu d'un seul coup d'œil.
  \item \textbf{Système de classification à trois niveaux}:
  \begin{itemize}
    \item \textbf{Niveau 1}: par type de fichier (document / image / vidéo / audio)
    \item \textbf{Niveau 2}: par scénario d'usage (travail / vie / apprentissage)
    \item \textbf{Niveau 3}: par événement ou projet spécifique (ex. « 2026Q2\_Projet A »)
  \end{itemize}
  \item \textbf{Plan de nettoyage mensuel}: fixez une date chaque mois (idéalement le jour de paie ou le 1er du mois), faites un tri ciblé du répertoire « Téléchargements » :\textbf{archivez l'utile, supprimez l'inutile, sauvegardez le précieux}. 12 petits rangements par an sont bien plus légers qu'un grand ménage annuel.
\end{itemize}

\begin{tcbpractice}{Liste des tâches du mois}
Choisissez un week-end libre, procédez dans l'ordre suivant :

\begin{enumerate}
  \item Créez dans le gestionnaire de fichiers un dossier professionnel du mois au format\texttt{2026Q2\_Projet A};
  \item Passez en revue tous les fichiers de la semaine dans le répertoire « Téléchargements », gardez l'utile, archivez, supprimez l'inutile ;
  \item Créez des dossiers annuels pour les photos familiales de cette année au format\texttt{Famille 2025}, le\texttt{Famille 2026}.
\end{enumerate}

\end{tcbpractice}

\booksubsection{Stratégie de traitement des fichiers redondants}

Avec le temps, le smartphone « grossit », souvent à cause des fichiers redondants.

\begin{itemize}
  \item \textbf{Nettoyage automatique 72 heures}: les fichiers temporaires reçus via les apps sociales telles que WeChat et QQ (stickers, images de discussion, courtes vidéos) sont généralement inutiles après un coup d'œil. Dans [WeChat] > [Réglages] > [Général] > [Espace de stockage], activez « nettoyage automatique », réglez le cache expiré à 72 heures. Le smartphone est instantanément plus léger, sans perdre de records importants.
  \item \textbf{Nettoyage ciblé de l'espace de stockage}: Android [Réglages] > [Stockage] > [Nettoyage accéléré], iOS [Réglages] > [Général] > [Espace iPhone], le système liste par catégories « cache / fichiers temporaires / apps désinstallables / gros fichiers ». Ciblez : (1) fichiers de cache ; (2) paquets d'installation expirés (anciens paquets d'installation laissés après la mise à jour des apps) ; (3) fichiers en double (photos de même nom et même taille).
  \item \textbf{Règles de classification automatique}: la plupart des apps d'album de marque prennent en charge la classification intelligente : personnes, paysages, documents, captures d'écran. Une fois activée, une nouvelle photo de groupe va automatiquement dans « Personnes > Zhang San », un contrat scanné va dans « Documents ».\textbf{Combiné au rangement mensuel, cela réduit considérablement la probabilité de « ne pas trouver »}.
\end{itemize}

\booksubsection{Norme de traitement sécurisé des images}

Les images sont le type de données le plus « sensible » du smartphone — une photo prise au hasard peut révéler votre position, votre appareil, l'heure. Nous recommandons\textbf{d'établir un flux de travail « double protection »}:

\begin{tcbwarn}{Procédure standard pour l'envoi de photos d'identité / captures sensibles}
\begin{enumerate}
  \item \textbf{Masquer les informations visibles}: avec l'outil intégré au système ou un outil tiers (comme Meitu XiuXiu, Snapseed)\textbf{barrez quelques chiffres du milieu du numéro de carte d'identité, la plaque d'immatriculation, le numéro de porte, le nom, le montant}et autres champs clés ;
  \item \textbf{Suppression des métadonnées}: avant l'envoi, dans l'album, cliquez sur « Partager » ou « Exporter », le système demande s'il faut conserver les informations de localisation, choisissez\textbf{« Ne pas conserver »}. Ainsi, le modèle d'appareil photo, les coordonnées GPS, l'heure de prise de vue ne suivront pas l'image lors de l'envoi ;
  \item \textbf{Filigrane temporaire}(optionnel) : ajoutez sur la capture sensible un filigrane textuel « Uniquement pour usage XX », pour éviter tout abus.
\end{enumerate}

\end{tcbwarn}

Méthodes courantes de nettoyage des métadonnées EXIF (opérations de base ; pour une protection approfondie, voir le chapitre 6 « Désensibilisation des photos ») :

\begin{itemize}
  \item \textbf{Intégré au système}: iOS [Réglages] > [Confidentialité et sécurité] > [Appareil photo], permet de désactiver « Position précise » ; lors du partage, choisir « Uniquement cette photo »\textbf{supprime automatiquement la plupart des EXIF}.
  \item \textbf{Android}: la plupart des apps d'album de marque proposent « Supprimer les informations de localisation » lors du partage, cochez.
  \item \textbf{Outils professionnels}: ExifTool (ordinateur), Metapho (iOS), Photo Metadata Remover (Android) peuvent supprimer complètement toutes les métadonnées.
\end{itemize}

Pour les images sensibles à\textbf{conserver à long terme}(carte d'identité, carte bancaire, scans de livret de famille, passeport), il est recommandé de les stocker dans un\textbf{album chiffré}:

\begin{itemize}
  \item iPhone : [Photos] > sélectionnez la photo > Masquer (visible uniquement après authentification Face ID) ;
  \item Huawei/Xiaomi : [Album] > appui long sur la photo > « Déplacer vers le coffre-fort » ;
  \item Apps de chiffrement tierces :\textbf{360 Album sécurisée}, le\textbf{Gestionnaire d'album Tencent (album privé)}, le\textbf{KeepSafe}, etc., protection par mot de passe distinct, discrétion accrue.
\end{itemize}
